\documentclass[12pt,a4paper]{article}

%% PACHETE MATEMATICE
\usepackage{amsmath,amssymb,amsfonts}
\usepackage{amsthm}
\usepackage{mathpazo}
\usepackage{eulervm}
\usepackage{enumerate}
\usepackage{color}
\usepackage{graphicx}
\usepackage[all]{xy}
\usepackage{xcolor}
\usepackage{framed}
\usepackage[unicode]{hyperref}
\setcounter{MaxMatrixCols}{20}
\usepackage{multicol}
% \usepackage[romanian]{babel}


%% ASPECT
\linespread{1.05}
\usepackage[T1]{fontenc}
\usepackage[utf8x]{inputenc}
\usepackage[margin=2cm]{geometry}
% \usepackage[color]{showkeys}
\usepackage{anyfontsize}
\usepackage{titlesec}
\titleformat*{\section}{\large\bfseries}

\usepackage{fancyhdr}
\pagestyle{fancy}
\rhead{\small \color{gray} Notițe de Adrian Manea}
\lhead{\small \color{gray} ETTI-AM2, 2017---2018, M. Joița \& A. Niță}


\usepackage[autostyle = false, style = german]{csquotes}
\usepackage[romanian]{babel}
\newcommand\qq{\enquote}

%% COMENZILE MELE
\renewcommand*{\proofname}{Dem.:}
\newcommand\bb{\mathbb}
\newcommand\dr{\mathrm}
\newcommand\kal{\mathcal}
\newcommand\fk{\mathfrak}
\newcommand\op{\oplus}
\newcommand\ot{\otimes}
\newcommand\ds{\displaystyle}
\newcommand\id{\indent}
\newcommand\nid{\noindent}
\newcommand\seq{\subseteq}
\newcommand\sneq{\subsetneq}
\newcommand\speq{\supseteq}
\newcommand\spneq{\supsetneq}
\newcommand\rar{\rightarrow}
\newcommand\Rar{\Rightarrow}
\newcommand\xrar{\xrightarrow}
\newcommand\lar{\leftarrow}
\newcommand\Lar{\Leftarrow}
\newcommand\xlar{\xleftarrow}
\newcommand\lrar{\leftrightarrow}
\newcommand\Lrar{\Leftrightarrow}
\newcommand\ideal{\trianglelefteq}
\newcommand\ai{\text{ a.\^{i}. }}
\newcommand\sk{(\textit{Schi\c{t}\u{a}}) }
\newcommand\rhup{\rightharpoonup}
\newcommand\rhdn{\rightharpoondown}
\newcommand\lhup{\leftharpoonup}
\newcommand\lhdn{\leftharpoondown}
\newcommand\vid{\emptyset}
\newcommand\wh{\widehat}
\newcommand\ol{\overline}
\newcommand\NN{\mathbb{N}}
\newcommand\RR{\mathbb{R}}
\newcommand\ZZ{\mathbb{Z}}
\newcommand\QQ{\mathbb{Q}}
\newcommand\CC{\mathbb{C}}
\newcommand\Lap{\mathcal{L}}
\newcommand\Rez{\mathrm{Rez}}

%% STILURILE MELE
\newtheoremstyle{dotless}{}{}{\itshape}{}{\bfseries}{:}{ }{}

\theoremstyle{dotless}
\newtheorem{theorem}{Teorem\u{a}}
\newtheorem{proposition}{Propozi\c{t}ie}
\newtheorem{lemma}{Lem\u{a}}
\newtheorem{corollary}{Corolar}
\newtheorem{exercise}{Exerci\c{t}iu}

\newtheoremstyle{dotlessdr}{}{}{}{}{\bfseries}{:}{ }{}
\theoremstyle{dotlessdr}
\newtheorem{example}{Exemplu}
\newtheorem{definition}{Defini\c{t}ie}
\newtheorem{remark}{Observa\c{t}ie}

\newenvironment{solution}
{\begin{proof}[Soluție:\nopunct]}
  {\end{proof}}


\begin{document}
% \definecolor{labelkey}{RGB}{222,0,0}

\begin{center}
  {\large\textbf{Seminar 11 \\
      Transformarea Laplace --- Aplicații \\
      Transformarea $ Z $}}
\end{center}

\vspace{2cm}

\section*{Ecuații și sisteme diferențiale}

\indent\indent Folosind transformata Laplace, putem rezolva ecuații și sisteme
diferențiale. Cu ajutorul proprietăților transformatei Laplace, aceste ecuații
și sisteme devin \emph{sisteme algebrice}, pe care le putem rezolva mult mai
simplu.

Principalele avantaje ale rezolvării ecuațiilor și sistemelor diferențiale cu
ajutorul transformatei Laplace sînt:
\begin{itemize}
\item rezolvarea unei ecuații neomogene se face \emph{direct}, nu mai este necesar
  să se trateze cazul omogen mai întîi;
\item valorile inițiale sînt folosite direct în calcul, nu se mai obține o soluție
  pînă la o constantă, care se determină din valorile inițiale.
\end{itemize}

Pentru aceasta, avem nevoie de o proprietate a transformatei Laplace aplicată
derivatei originalului. Folosind definiția și integrarea prin părți, se pot verifica
simplu următoarele proprietăți:
\begin{align*}
  \Lap (f') &= s \Lap (f) - f(0) \\
  \Lap (f^{\prime\prime}) &= s^2 \Lap (f) - sf(0) - f'(0).
\end{align*}

De fapt, în general, avem:
\[
  \Lap (f^{(n)}) = s^n \Lap (f) - s^{n-1} f(0) - s^{n-2} f'(0) - \dots - %
  f^{(n-1)}(0).
\]

De asemenea, pentru integrale, avem deja teorema integrării originalului:
\[
  \Lap \int_0^t f(\tau) d\tau = \frac{1}{s} F(s), \quad \text{unde } %
  F(s) = \Lap f(t).
\]

Rezultă, folosind transformarea inversă:
\[
  \int_0^t f(\tau) d\tau = \Lap^{-1} \Big( \frac{1}{s} F(s) \Big).
\]

\vspace{1cm}

De exemplu, pentru a rezolva ecuația diferențială:
\[
  y^{\prime \prime} + ay^\prime + by = r(t), \quad y(0) = K_0, y'(0) = K_1,
\]
aplicăm transformata Laplace și folosim proprietățile de mai sus. Fie $ Y = \Lap y(t) $

Se obține ecuația algebrică:
\[
  (s^2 Y - sy(0) - y'(0)) + a(sY - y(0)) + bY = R(s),
\]
unde $ R(s) = \Lap r $. Forma echivalentă este:
\[
  (s^2 + as + b) Y = (s + a)y(0) + y'(0) + R(s).
\]

Împărțim prin $ s^2 + as + b $ și folosim formula:
\[
  Q(s) = \frac{1}{s^2 + as + b} =
  \frac{1}{(s + \frac{1}{2}a)^2 + b - \frac{1}{4}a^2},
\]
de unde rezultă:
\[
  Y(s) = \big( (s + a) y(0) + y'(0)\big) Q(s) + R(s) Q(s).
\]

În forma aceasta, descompunem $ Y(s) $ în fracții simple, dacă este nevoie și
folosim tabelul de transformate Laplace, pentru a afla $ y = \Lap^{-1}(Y) $.

\vspace{1cm}

De exemplu:
\[
  y^{\prime\prime} - y = t, \quad y(0) = 1, y'(0) = 1.
\]

\emph{Soluție:} Aplicăm transformata Laplace și ajungem la ecuația:
\begin{align*}
  s^2 Y - sy(0) - y'(0) - Y &= \frac{1}{s^2} \\
  (s^2 - 1)Y &= s + 1 + \frac{1}{s^2}.
\end{align*}

Rezultă $ Q = \frac{1}{s^2 - 1} $ și ecuația devine:
\begin{align*}
  Y &= (s + 1)Q + \frac{1}{s^2} Q \\
    &= \frac{s + 1}{s^2 - 1} + \frac{1}{s^2(s^2 - 1)} \\
    &= \frac{1}{s - 1} + \Big( \frac{1}{s^2 - 1} - \frac{1}{s^2} \Big)
\end{align*}

Folosind tabelul și proprietățile transformatei Laplace, obținem soluția:
\begin{align*}
  y(t) &= \Lap^{-1} Y \\
       &= \Lap^{-1} \Big( \frac{1}{s^{-1}} \Big) + \Lap^{-1} \Big( \frac{1}{s^2 - 1} \Big) - %
         \Lap^{-1} \Big( \frac{1}{s^2} \Big) \\
       &= e^t + \sinh t - t.
\end{align*}

\vspace{1cm}

%%%%%%%%%%%%%%%%%%%%%%%%%%%%%%%%%%%%%%%%%%%%%%%%%%%%%%%%%%%%%%%%%%%%%% 

\section*{Exerciții}

\indent\indent 1. Să se rezolve următoarele probleme Cauchy, folosind transformata
Laplace:
\begin{enumerate}[(a)]
\item $ y'(t) + 2y(t) = 4t, y(0) = 1 $ ;
\item $ y'(t) + y(t) = \sin 4t, y(0) = 0 $;
\item $ y^{\prime\prime}(t) + 2y'(t) + 5y(t) = 0, y(0) = 0, y'(0) = 0 $;
\item $ 2y^{\prime\prime}(t) - 6y'(t) + 4y(t) = 3e^{3t}, y(0) = 1, y'(0) = -1 $;
\item $ y^{\prime\prime}(t) - 2y(t) + y(t) = e^t, y(0) = -2, y'(0) = -3 $;
\item $ y^{\prime\prime}(t) + 4y(t) = 3\cos^2(t), y(0) = 1, y'(0) = 2 $.  
\end{enumerate}


\vspace{1cm}

2. Să se rezolve următoarele sisteme diferențiale:
\begin{enumerate}[(a)]
\item $
  \begin{cases}
    x' + x + 4y &= 10 \\
    x - y' - y &= 0 \\
    x(0) = 4, & y(0) = 3
  \end{cases}
  $
\item $
  \begin{cases}
    x' + x - y &= e^t \\
    y' + y - x &= e^t \\
    x(0) = 1, & y(0) = 1
  \end{cases}
  $
\item $
  \begin{cases}
    x' + 2y' + x - y &= 5 \sin t \\
    2x' + 3y' + x - y &= e^t \\
    x(0) = 2, & y(0) = 1
  \end{cases}
  $
\end{enumerate}

\emph{Indicație:} Aplicăm transformata Laplace fiecărei ecuații
și notăm $ \Lap x(t) = X(s) $ și $ \Lap y(t) = Y(s) $. Apoi rezolvăm sistemul
\emph{algebric} obținut cu necunoscutele $ X $ și $ Y $, cărora la final le
aplicăm transformata Laplace inversă.


%%%%%%%%%%%%%%%%%%%%%%%%%%%%%%%%%%%%%%%%%%%%%%%%%%%%%%%%%%%%%%%%%%%%%%

\section*{Transformata $ Z $}

\indent\indent Definim acum o nouă transformată, dar de data aceasta pe un
\emph{caz discret}.

\begin{definition}\label{def:semnal-discret}
  Se numește \emph{semnal discret} o funcție $ x : \ZZ \to \CC $, dată de
  $ n \mapsto x_n $ (sau, echivalent, $ x(n) $ ori $ x[n] $).

  Mulțimea semnalelor discrete se va nota cu $ S_d $.

  Dacă $ x_n = 0 $ pentru $ n < 0 $, spunem că semnalul are \emph{suport pozitiv},
  iar mulțimea lor se va nota cu $ S_d^+ $.
\end{definition}

Un semnal particular este următorul: pentru $ k \in \ZZ $ fixat, definim:
\[
  \delta_k(n) =
  \begin{cases}
    1, & n = k \\
    0, & n \neq k
  \end{cases}.
\]
Acesta se numește \emph{impulsul unitar discret} la momentul $ k $ și vom nota
$ \delta_0 $ cu $ \delta $ simplu.

Cîteva noțiuni și operații specifice cu semnale urmează.

\begin{definition}\label{def:intirziat}
  Fie $ x \in S_d $ și $ k \in \ZZ $ fixat arbitrar. Semnalul $ y = (x_{n-k})_{n \in \ZZ} $
  se numește \emph{întîrziatul} lui $ x $ cu $ k $ momente.
\end{definition}

Operația de \emph{convoluție} o reîntîlnim și în cazul semnalelor discrete:

\begin{definition}\label{def:convolutie}
  Fie $ x, y \in S_d $. Dacă seria:
  \[
    \sum_{k \in \ZZ} x_{n - k} y_k
  \]
  este convergentă pentru orice $ n \in \ZZ $ și are suma $ z_n $, atunci semnalul
  $ z = (z_n)_n $ se numește \emph{convoluția} semnalelor $ x $ și $ y $ și se notează
  $ z = x \ast y $.
\end{definition}

Să remarcăm că, din definiție, avem pentru $ x, y \in S_d^+ $ existența $ x \ast y $
și $ x \ast y = y \ast x $. În plus, au loc:
\[
  x \ast \delta = x \quad \text{și} \quad (x \ast \delta_k)(n) = x_{n - k}.
\]

Ajungem în fine la definiția principală:
\begin{definition}\label{def:transf-z}
  Fie $ s \in S_d $, cu $ s = (a_n)_n $. Se numește \emph{transformata $ Z $} sau
  \emph{transformata Laplace discretă} a acestui semnal funcția complexă definită
  prin:
  \[
    L_s(z) = \sum_{n \in \ZZ} a_n z^{-n},
  \]
  care se definește în domeniul de convergență al seriei Laurent corespunzătoare.
\end{definition}

Principalele proprietăți ale transformării $ Z $ sînt:

\begin{enumerate}[(1)]
\item Există $ R, r > 0 $ astfel încît seria care definește transformarea $ Z $ să fie
  convergentă în coroana $ r < |z| < R $;
\item \textbf{Liniaritatea:} Asocierea $ s \mapsto L_s $ este $ \CC $-liniară și injectivă,
  deci:
  \[
    L_{\alpha_1 s_1 + \alpha_2 s_2}(z) = \alpha_1 L_{s_1}(z) + \alpha_2 L_{s_2}(z), %
    \forall \alpha_1, \alpha_2 \in \CC, s_1, s_2 \in S_d.
  \]
\item Dacă $ s \in S_d^+ $, cu $ s = (a_n) $, atunci $ \ds\lim_{z \to \infty} L_s(z) = a_0 $,
  iar dacă există limita $ \ds\lim_{n \to \infty} a_n = \ell $, atunci:
  \[
    \lim_{z \to 1} \frac{z - 1}{z} L_s(z) = 1.
  \]
\item \textbf{Inversa transformării $ Z $:} Fie $ s \in S_d^+ $, cu $ s = (a_n) $. Presupunem
  că funcția $ L_s(z) $ este olomorfă în domeniul $ r < |z| < R $. Pentru orice $ r < \rho < R $,
  fie $ \gamma_\rho $ frontiera discului $ |z| \leq \rho $ parcursă în sens pozitiv o singură
  dată. Atunci avem:
  \[
    a_n = \frac{1}{2 \pi i} \int_{\gamma_\rho} z^{n-1} L_s(z) dz, \quad n \in \ZZ.
  \]
\item \textbf{Teorema de convoluție:} Fie $ s, t \in S_d^+ $. Atunci $ s \ast t \in S_d^+ $ și
  are loc:
  \[
    L_{s \ast t} = L_s \cdot L_t.
  \]
  În particular:
  \[
    L_{s \ast \delta_k}(z) = z^{-k} L_s(z), \quad k \in \ZZ.
  \]
\item \textbf{Prima teoremă de întîrziere:} Pentru $ n \in \NN^\ast $:
  \[
    \Lap_s (f(t - n)) = z^{-n} \Lap_s(f).
  \]
\item \textbf{A doua teoremă de întîrziere (teorema de deplasare):}
  \[
    \Lap_s (f (t + n)) = z^n \cdot \Big( \Lap_s(f) - \sum_{t = 0}^{n-1} f(t) z^{-t} \Big), \quad %
    n \in \NN^\ast.
  \]
\end{enumerate}

În tabelul din figura~\ref{fig:transf-z} sînt prezentate transformatele $ Z $ ale funcțiilor
uzuale.

\begin{figure}[!htbp]
  \centering
  \begin{tabular}{c | c}
    $ s $ & $ L_s $ \\
    \hline\hline
    $ \begin{cases}
      h_n = 0, & n < 0 \\
      h_n = 1, & n \geq 0
    \end{cases} $ & $ \dfrac{z}{z - 1} $ \\
          & \\
    $ \delta_k, k \in \ZZ $ & $ \dfrac{1}{z^k} $ \\
          &\\
    $ s = (n)_{n \in \NN} $ & $ \dfrac{z}{(z - 1)^2} $ \\
          &\\
    $ s = (n^2)_{n \in \NN} $ & $ \dfrac{z(z+1)}{(z - 1)^3} $ \\
          &\\
    $ s = (a^n)_{n \in \NN}, a \in \CC $ & $ \dfrac{z}{z - a} $ \\
          &\\
    $ s = (e^{an})_{n \in \NN}, a \in \RR $ & $ \dfrac{z}{z - e^a} $ \\
          &\\
    $ s = (\sin(\omega n))_{n \in \NN}, \omega \in \RR $ & $ \dfrac{z \sin \omega}{z^2 - 2z \cos \omega + 1} $ \\
          &\\
    $ s = (\cos(\omega n))_{n \in \NN}, \omega \in \RR $ & $ \dfrac{z(z - \cos \omega)}{z^2 - 2z \cos \omega + 1} $
  \end{tabular}
  \caption{\emph{Transformate $ Z $ uzuale}}\label{fig:transf-z}
\end{figure}

%%%%%%%%%%%%%%%%%%%%%%%%%%%%%%%%%%%%%%%%%%%%%%%%%%%%%%%%%%%%%%%%%%%%%%

\section*{Exerciții}

\indent\indent 1. Să se determine semnalul $x \in S_d^+$, a cărui transformată $Z$ este 
dată de:
\begin{enumerate}[(a)]
\item $\Lap_s (z) = \dfrac{z}{(z-3)^2}$ ;
\item $\Lap_s (z) = \dfrac{z}{(z-1)(z^2 + 1)}$ ;
\item $\Lap_s (z) = \dfrac{z}{(z-1)^2(z^2 + z - 6)}$ ;
\item $\Lap_s (z) = \dfrac{z}{z^2 + 2az + 2a^2}, a > 0$ parametru.
\end{enumerate}

\begin{solution}
(a) Avem:
\begin{align*}
x_n &= \dfrac{1}{2\pi i} \int_{|z| = \rho} z^{n-1} \Lap_s (z) dz \\
    &= \Rez(z^{n-1} \Lap_s(z), 3) \\
    &= \Rez \Big( \dfrac{z^n}{(z-3)^2}, 3 \Big) \\
    &= \lim_{z \to 3} \Big( (z-3)^2 \cdot \dfrac{z^n}{(z-3)^2} \Big)^{\prime} \\
    &= \lim_{z \to 3} nz^{n-1} \\
    &= n3^{n-1}.
\end{align*}

(b)
\begin{align*}
x_n &= \dfrac{1}{2 \pi i} \int_{|z| = \rho} z^{n-1} \Lap_s(z) dz \\
    &= \Rez(z^{n-1} \Lap_s (z), 1) + \Rez(z^{n-1} \Lap_s(z), i) + 
    \Rez(z^{n-1} \Lap_s(z), -i) \\
\Rez(z^{n-1}\Lap_s(z),1) &= \lim_{z \to 1} z^{n-1} \dfrac{z}{(z - 1)(z^2 + 1)} 
\cdot (z-1) = \dfrac{1}{2} \\
\Rez(z^{n-1} \Lap_s(z), i) &= \dfrac{i^n}{2i \cdot (i-1)} \\
\Rez(z^{n-1} \Lap_s(z), -i) &= \dfrac{(-1)^n i^n}{2i(i+1)}.
\end{align*}

Remarcăm că pentru $n = 4k$ și $n = 4k+1$, avem $x_n = 0$, iar în celelalte 
două cazuri, $x_n = 1$. \\

(c)
\begin{align*}
\Rez(z^{n-1} \Lap_s(z), 1) &= \lim_{z\to 1} \Big[ z^{n-1} \cdot (z-1)^2 \cdot 
\dfrac{z^2}{(z-1)^2 \cdot (z^2 + z - 6)} \Big]^{\prime}\\
                       &= - \dfrac{4n + 3}{16}. \\
\Rez(z^{n-1} \Lap_s(z), 2) &= \dfrac{2^n}{5} \\
\Rez(z^{n-1} \Lap_s(z), -3) &= -\dfrac{(-3)^n}{80}.
\end{align*}

Obținem $x_n = - \dfrac{4n+3}{16} + \dfrac{2^n}{5} - \dfrac{(-3)^n}{80}$. \\

(d) $z_{1,2} = a (-1 \pm i)$ sînt poli simpli. Avem:
\begin{align*}
x_n &= \Rez\Big( \dfrac{z^n}{(z^2 + 2a + 2a^2)}, z_1 \Big) + 
\Rez\Big(  \dfrac{z^n}{(z^2 + 2a + 2a^2)}, z_2 \Big) \\
    &= \dfrac{a^n(-1 + i)^n}{2z_1 + 2a} + \dfrac{a^n(-1-i)^n}{2z_1 + 2a} \\
    &= -\dfrac{i}{2a}(z_1^n - z_2^n).
\end{align*}

Putem scrie trigonometric numerele $z_1$ și $z_2$:
\begin{align*}
z_1 &= a(-1 + i) = a \sqrt{2} \Big( \cos \dfrac{3\pi}{4} + 
i \sin \dfrac{3 \pi}{4} \Big) \\
z_2 &= a(-1 - i) = a \sqrt{2} \Big( \cos \dfrac{3 \pi}{4} - 
i \sin \dfrac{3 \pi}{4} \Big).
\end{align*}

Deci: $x_n = 2^{\frac{n}{2}} a^{n-1} \sin \dfrac{3 n \pi}{4}$.
\end{solution}

%%%%%%%%%%%%%%%%%%%%%%%%%%%%%%%%%%%%%%%%%%%%%%%%%%%%%%%%%%%%%%%%%%%%%%%%%%%%%%%%

\vspace{1cm}

2. Fie $x = (x_n) \in S_d^+$ și $y = (y_n)$, unde $y_n = x_0 + \dots + x_n$. 
Să se arate că $Y(z) = \dfrac{z}{z-1} X(z)$.

\begin{solution}
Avem $Y(z) = \displaystyle\sum_{n=0}^\infty y_n z^{-n}.$ Dar: 
\[
X(z) = \sum_{n=0}^\infty x_n z^{-n} \text{ și } 
\sum_{n=0}^\infty x_{n-1} z^n = \dfrac{1}{z} X(z),
\]
deoarece $x_{-1} = 0$. Putem continua și obținem 
$\displaystyle\sum_{n=0}^\infty x_{n-k} z^{-n} = \dfrac{1}{z^k} X(z)$. Așadar: 
\[
Y(z) = X(z) \cdot \Big( 1 + \dfrac{1}{z} + \dfrac{1}{z^2} + \dots \Big) 
= X(z) \cdot \dfrac{z}{z-1}.
\]
\end{solution}

\vspace{1cm}

3. Cu ajutorul transformării $Z$, să se determine șirurile $(x_n)$ definite 
prin următoarele relații:
\begin{enumerate}[(a)]
\item $x_0 = 0, x_1 = 1, x_{n+2} = x_{n+1} + x_n, n \in \mathbb{N}$ 
(șirul lui Fibonacci);
\item $x_0 = 0, x_1 = 1, x_{n+2} = x_{n+1} - x_n, n \in \mathbb{N}$;
\item $x_0 = x_1 = 0, x_2 = -1, x_3 = 0, x_{n+4} + 
2x_{n+3} + 3 x_{n+2} + 2 x_{n+1} + x_n = 0, n \in \mathbb{N}$;
\item $x_0 = 2, x_{n+1} + 3x_n = 1, n \in \mathbb{Z}$;
\item $x_0 = 0, x_1 = 1, x_{n+2} - 4 x_{n+1} + 3x_n 
= (n+1)4^n, n \in \mathbb{N}$.
\end{enumerate}

\begin{solution}
Abordarea generală este să considerăm șirul $(x_n)$ ca fiind restricția unui 
semnal $x \in S_d^+$ la $\mathbb{N}$ și rescriem relațiile de recurență sub 
forma unor ecuații de convoluție $a \ast x = y$, pe care le rezolvăm în $S_d^+$.\\

(a) Fie $x \in S_d^+$, astfel încît restricția lui la $\mathbb{N}$ să fie șirul 
căutat. Deoarece avem: 
\[
x_{n+2} - x_{n+1} - x_n = y_n, n \in \mathbb{Z},
\]
cu $y_n = 0$ pentru $n \neq -1$ și $y_{-1} = 1$, avem ecuația de convoluție: 
\[
a \ast x = y, \text{ unde } a 
= \delta_{-2} + \delta_{-1} + \delta, y = \delta_{-1}.
\]

Aplicăm transformata $Z$ și rezultă: 
\[
\Lap_s x(z) (z^2 - z - 1) = z \Longrightarrow x_n 
= \dfrac{1}{\sqrt{5}} \Big[ \Big( \dfrac{1 + \sqrt{5}}{2} \Big)^n - 
\Big( \dfrac{1 - \sqrt{5}}{2} \Big)^n \Big].
\]

(b) Ca în cazul anterior, avem $a \ast x = y$, cu 
$a = \delta_{-2} - \delta_{-1} + \delta$, unde $y = \delta_{-1}$. Aplicînd 
transformata $Z$, obținem: 
\[
\Lap_s x(z) (z^2 - z + 1) = z \Longrightarrow \Lap_s x(z) 
= \dfrac{z}{z^2 - z + 1}.
\]

Obținem:
\[
x_n = \Rez \Big( z^{n-1} \dfrac{z}{z^2 - z + 1}, 
\dfrac{1 + i \sqrt{3}}{2} \Big) + \Rez \Big( z^{n-1} \dfrac{z}{z^2 - z + 1}, 
\dfrac{1 - i \sqrt{3}}{2} \Big).
\]

Calculăm reziduurile, cu notația $\varepsilon = \dfrac{1 + i \sqrt{3}}{2}$ și 
$\overline{\varepsilon} = \dfrac{1 - i \sqrt{3}}{2}$:
\begin{align*}
\Rez \Big( \dfrac{z^n}{z^2 - z + 1}, \varepsilon \Big) 
&= \lim_{z \to \varepsilon} \dfrac{z^n}{z^2 - z + 1} 
\big( z - \varepsilon \big) \\
&= \dfrac{\varepsilon^n}{i \sqrt{3}} \\
&= \dfrac{\cos \frac{2 n \pi}{3} + i \sin \frac{2 n \pi}{3}}{i \sqrt{3}}.
\end{align*}

Similar: 
\[
\Rez \Big( \dfrac{z^n}{z^2 - z + 1}, \overline{\varepsilon} \Big) 
= \dfrac{\cos \frac{2n\pi}{3} - i \sin \frac{2 n \pi}{3}}{-i\sqrt{3}}.
\]

Rezultă: 
\[
x_n = \dfrac{2}{\sqrt{3}} \sin \dfrac{2n\pi}{3}, n \in \mathbb{N}.
\]

(c) Ecuația $a \ast x = y$ este valabilă pentru: 
\[
a = \delta_{-4} + 2 \delta_{-3} + 3 \delta_{-2} + 2 \delta_{-1} + \delta, 
\quad y = -\delta_{-2} - 2 \delta_{-1}.
\]

Aplicăm transformata $Z$ și obținem: 
$\Lap_s x(z) = - \dfrac{z(z+2)}{(z^2 + z + 1)^2}$. Descompunem în fracții 
simple, calculăm reziduurile și ținem cont de faptul că rădăcinile numitorului, 
$\varepsilon_1, \varepsilon_2$ sînt poli de ordinul 2, obținem: 
\[
x_n = \dfrac{(2n-4)(\varepsilon_1^n - \varepsilon_2^n) - 
(n+1)(\varepsilon_1^{n-1} + \varepsilon_1^{n-2} - \varepsilon_2^{n-1} - 
\varepsilon_2^{n-2})}{(\varepsilon_2 - \varepsilon_1)^3} 
= \dfrac{2(n-1)}{\sqrt{3}} \sin \dfrac{2 n \pi}{3}, n \in \mathbb{N}.
\]

(d) Ecuația corespunzătoare este $a \ast x = y$, cu $a = \delta_{-1} + 3 \delta$ și 
$y_n = 1, \forall n \geq 1$, iar $y_{-1} = x_0 + 3 x_{-1} = 2$, cu 
$y_n = 0, \forall n \leq -2$, adică $y = 1 + 2 \delta_{-1}$.

Așadar: 
\[
\delta_{-1} \ast x + 3 \delta \ast x = 1 + 2 \delta_{-1}.
\]

Aplicăm transformata $Z$ și obținem:
\begin{align*}
z \Lap_s x(z) + 3 \Lap_s x(z) &= \dfrac{z}{z-1} + 2z \\
                              &= \dfrac{2z^2 + 3z}{z-1} \\
\Longrightarrow \Lap_s x(z) &= \dfrac{2z^2 + 3z}{(z-1)(z+3)} \\
\Longrightarrow x_n &= \Rez(z^{n-1} \cdot \dfrac{2z^2 + 3z}{(z-1)(z+3)}, 1) + 
\Rez (z^{n-1} \cdot \dfrac{2z^2 + 3z}{(z-1)(z+3)}, -3) \\
\Rez(z^{n-1} \cdot \dfrac{2z^2 + 3z}{(z-1)(z+3)}, 1) &= \lim_{z \to 1} (z-1) 
\cdot z^{n-1} \cdot \dfrac{2z^2 + 3z}{(z-1)(z+3)} = \dfrac{5}{4} \\
\Rez(z^{n-1} \cdot \dfrac{2z^2 + 3z}{(z-1)(z+3)}, -3) 
&= \lim_{z \to -3} (z+3) z^{n-1} \cdot \dfrac{2z^2 + 3z}{(z-1)(z+3)} \\
&= (-3)^{n-1} \cdot \dfrac{12}{-4} = -3 \cdot (-3)^{n-1}.	                              
\end{align*}

Rezultă: $x_n = \dfrac{5}{4} - 3 \cdot (-3)^{n-1}$. \\

(e) Avem ecuația: $a \ast x = y$, unde 
$a = \delta_{-2} - 4 \delta_{-1} + 3 \delta$, cu 
$y_n = 0, \forall n \leq -2, y_{-1} = 1$ și 
$y_n = (n+1)4^n, \forall n \in \mathbb{N}$.

Fie $s_1 = (n4^n)_n, s_2 = (4^n)_n$. Atunci:
\begin{align*}
\Lap_s s_1(z) &= -z \Lap_s s_2^{\prime}(z) 
= - z \big( \dfrac{z}{z-4} \big)^{\prime} = \dfrac{4z}{(z-4)^2} \\
\Longrightarrow \Lap_s x(z)(z^2 - 4z + 3) 
&= \dfrac{4z}{(z-4)^2} + \dfrac{z}{z-4} + z \\
&= \dfrac{z^2}{(z-4)^2} + z \\
\Longrightarrow  \Lap_s x(z) &= \dfrac{z(z^2 - 7z + 16)}{(z-4)^2 (z-1) (z-3)}.
\end{align*}

Descompunem în fracții simple și obținem, în fine: 
\[
x_n = \dfrac{1}{9} \big[ 18 \cdot 3^n + (3n - 13) 4^n - 5 \big], 
n \in \mathbb{N}.
\]
\end{solution}

\textbf{OBSERVAȚIE:} Toate exercițiile cu recurențe se mai pot rezolva în alte
două moduri:

(1) Se poate aplica teorema de convoluție relației de recurență. De exemplu,
din recurența:
\[
  x_{n + 2} - 2x_{n + 1} + x_n = 2
\]
putem obține:
\[
  Z(x_{n+2}) - 2Z(x_{n + 1}) + Z(x_n) = 2Z(1),
\]
iar $ Z(x_{n+2}) = Z(x_n \ast \delta_{-2}) = Z(x_n) \cdot Z(\delta_{-2}) $ etc.

(2) Se poate aplica teorema de deplasare. În aceeași recurență de mai sus,
de exemplu, avem:
\[
  Z(x_{n+2}) = z^n \Big( Z(x_n) - x_0 - x_1 z^{-1} \Big)
\]
și la fel pentru celelalte.

\end{document}
